% AA-26 | Iteracja XIII; szczegolowy audyt historyczny 2026-11-20, partia 02.
\chapter[Rozszerzenia cykliczne i teoria Kummera]{Rozszerzenia cykliczne, pierwiastki z jedynki i teoria Kummera}
\label{aa:ch:kummer}

\section{Motywacja: od pierwiastnika do grupy Galois}
Niech $F$ będzie ciałem, $n\geq2$ liczbą całkowitą, a $a\in F^\times$.
Po dołączeniu pierwiastka $\alpha$ równania $\alpha^n=a$
otrzymujemy rozszerzenie $E=F(\alpha)$. Chcielibyśmy odczytać
jego grupę Galois z prostego równania $X^n-a=0$.
Dwa zjawiska mogą jednak temu przeszkadzać: inne pierwiastki
tego wielomianu mogą nie należeć do $E$, a w dodatniej
charakterystyce wielomian może mieć pierwiastki wielokrotne.
Teoria Kummera określa warunki, w których oba problemy znikają,
i wyjaśnia, dlaczego powstają rozszerzenia cykliczne.

Dla ciała $F$ symbol $F^\times$ oznacza jego grupę
multiplikatywną elementów niezerowych. Przez
$\operatorname{char}F\nmid n$ rozumiemy, że obraz liczby całkowitej
$n$ w $F$ jest niezerowy, a więc odwracalny.
Wszystkie pierwiastki wielomianów rozpatrujemy w ustalonym
domknięciu algebraicznym $\overline F$.

\section{Pierwiastki z jedynki i rozszerzenia radykalne}
\begin{definicja}[Grupa pierwiastków z jedynki]\label{aa:def:mu-n}
Niech $F$ będzie ciałem i $n\geq1$ liczbą całkowitą. Określamy
\[
\mu_n(F)=\{u\in F^\times:u^n=1\}.
\]
Jest to podgrupa $F^\times$: $1\in\mu_n(F)$, iloczyn
dwóch elementów spełniających równanie $u^n=1$
spełnia je ponownie, a $u^{-1}$ także jest jego rozwiązaniem.
Element $\zeta\in\mu_n(F)$ nazywamy
\emph{pierwotnym $n$-tym pierwiastkiem z jedynki},
jeżeli jego rząd w $F^\times$ jest dokładnie równy $n$.
\end{definicja}

Najpierw uzasadnimy własność, którą często używa się bez dowodu:
skończona podgrupa grupy multiplikatywnej ciała jest cykliczna.

\begin{lemat}[Skończone podgrupy $F^\times$]
Dla dowolnego ciała $F$ każda skończona podgrupa $H\leq F^\times$
jest cykliczna.
\end{lemat}
\begin{proof}
Grupa $H$ jest abelowa, ponieważ mnożenie w ciele jest przemienne.
Jeśli $H=\{1\}$, teza jest spełniona. Załóżmy, że $|H|>1$
i oznaczmy przez $m$ najmniejszą wspólną wielokrotność
rzędów wszystkich elementów $H$. Jest ona określona,
ponieważ $H$ jest skończona.

Niech $p^e$ będzie dowolną maksymalną potęgą liczby pierwszej
$p$ dzielącą $m$. W definicji najmniejszej wspólnej wielokrotności
musi występować element $h_p\in H$, którego rząd dzieli się przez
$p^e$. Jeśli $\operatorname{ord}(h_p)=p^e q$ z $p\nmid q$,
to element $g_p=h_p^q$ ma rząd dokładnie $p^e$:
$(h_p^q)^{p^e}=1$, a żadna mniejsza dodatnia potęga $p$
nie wystarcza do uzyskania $1$.

Dla różnych dzielników pierwszych $p$ liczby $m$ elementy
$g_p$ komutują, a ich rzędy są parami względnie pierwsze.
W grupie abelowej iloczyn dwóch elementów o względnie
pierwszych rzędach $r,s$ ma rząd $rs$: jeśli
$(uv)^k=1$, to $u^k=v^{-k}$ należy jednocześnie
do podgrup $\langle u\rangle$ i $\langle v\rangle$; ich
przecięcie jest trywialne, bo rząd elementu przecięcia
dzieli jednocześnie $r$ i $s$. Stąd $r\mid k$ i $s\mid k$,
więc $rs\mid k$. Powtarzając ten argument, otrzymujemy,
że iloczyn $g=\prod_{p\mid m}g_p$ ma rząd
$\prod_{p\mid m}p^{v_p(m)}=m$.
W szczególności $H$ zawiera $m$ różnych potęg $g$.
Z drugiej strony rząd każdego elementu $h\in H$ dzieli $m$,
więc $h^m=1$. Wszystkie elementy $H$ są zatem pierwiastkami
wielomianu $X^m-1\in F[X]$. Ponieważ niezerowy wielomian
stopnia $m$ nad ciałem ma najwyżej $m$ różnych pierwiastków,
$|H|\leq m$. Istnienie elementu $g$ rzędu $m$ daje
$|H|\geq m$. Obie nierówności oznaczają
$H=\langle g\rangle$.
\end{proof}

\begin{wniosek}[Kiedy istnieje pierwotny pierwiastek?]
Jeśli $\operatorname{char}F\nmid n$ oraz wielomian $X^n-1$
ma w $F$ wszystkie swoje $n$ pierwiastków, to
$\mu_n(F)$ jest grupą cykliczną rzędu $n$.
\end{wniosek}
\begin{proof}
Pochodna $X^n-1$ jest równa $nX^{n-1}$.
Przy założeniu $\operatorname{char}F\nmid n$ wspólny pierwiastek
$X^n-1$ i $nX^{n-1}$ musiałby jednocześnie być zerem
$X$, co jest niemożliwe dla pierwiastka $X^n-1$.
Zatem wielomian ma $n$ różnych pierwiastków w $\overline F$.
Jeśli wszystkie należą do $F$, to $|\mu_n(F)|=n$.
Z poprzedniego lematu wynika cykliczność tej grupy,
a jej generator jest pierwotnym $n$-tym pierwiastkiem z jedynki.
\end{proof}

\begin{twierdzenie}[Kryterium cykliczności prostego rozszerzenia radykalnego]\label{aa:tw:kummer-radical-cyclic}
Niech $F$ będzie ciałem, $n\geq2$ liczbą całkowitą
odwracalną w $F$ i niech $\mu_n(F)$ ma rząd $n$.
Dla każdego $a\in F^\times$ i każdego $\alpha\in\overline F$
takiego, że $\alpha^n=a$, rozszerzenie $E=F(\alpha)/F$
jest skończonym rozszerzeniem Galois o grupie cyklicznej,
której rząd dzieli $n$.
\end{twierdzenie}
\begin{proof}
\emph{Rozkład wielomianu.}
Ponieważ $a\neq0$, także $\alpha\neq0$. Jeśli $\beta^n=a$,
to $(\beta/\alpha)^n=1$, zatem $\beta=\zeta\alpha$
dla pewnego $\zeta\in\mu_n(\overline F)$.
Na mocy założeń wszystkie $n$-te pierwiastki z jedynki
leżą w $F$, więc każdy pierwiastek $X^n-a$ należy do $E$.
Oznacza to, że $E$ jest ciałem rozkładu tego wielomianu nad $F$.

\emph{Rozdzielczość.}
Pochodna $X^n-a$ jest równa $nX^{n-1}$.
Ponieważ $n\neq0$ w $F$, wspólny pierwiastek obu wielomianów
musiałby być zerem $X$, ale $0^n-a=-a\neq0$.
Wszystkie pierwiastki są więc jednokrotne. Ponieważ $E$
jest ciałem rozkładu wielomianu rozdzielczego,
$E/F$ jest skończonym rozszerzeniem Galois.

\emph{Opis grupy.}
Niech $G=\operatorname{Gal}(E/F)$. Dla każdego $\tau\in G$
mamy $\tau(\alpha)^n=\tau(a)=a$, więc
$\chi(\tau):=\tau(\alpha)/\alpha\in\mu_n(F)$.
Sprawdźmy, że odwzorowanie $\chi:G\to\mu_n(F)$ jest
homomorfizmem, a nie tylko dowolną funkcją.
Jeśli $\tau,\rho\in G$, to $\rho(\alpha)=\chi(\rho)\alpha$,
a $\tau$ ustala element $\chi(\rho)\in F$. Stąd
\[
\chi(\tau\rho)=\frac{\tau(\rho(\alpha))}{\alpha}
=\frac{\chi(\rho)\tau(\alpha)}{\alpha}
=\chi(\rho)\chi(\tau)=\chi(\tau)\chi(\rho),
\]
gdzie ostatnia równość wykorzystuje przemienność grupy
$\mu_n(F)$. Jeśli $\chi(\tau)=1$, to $\tau(\alpha)=\alpha$;
automorfizm $\tau$ ustala też $F$, a zatem każdy element
$F(\alpha)=E$. Wobec tego $\tau=\operatorname{id}_E$,
więc $\chi$ jest iniektywne.
Grupa $G$ jest zatem izomorficzna z podgrupą cyklicznej grupy
$\mu_n(F)$ rzędu $n$. Każda podgrupa grupy cyklicznej
jest cykliczna, a na mocy twierdzenia Lagrange'a jej rząd
dzieli $n$.
\end{proof}

\begin{uwaga}[Dlaczego sam warunek „nie jest $n$-tą potęgą” nie wystarcza?]
Dla złożonego $n$ wielomian $X^n-a$ może być rozkładalny,
mimo że $a$ nie jest $n$-tą potęgą w $F$.
Poprzednie twierdzenie gwarantuje wtedy jedynie, że stopień
rozszerzenia \emph{dzieli} $n$, nie zaś, że jest równy $n$.
Przykład o stopniu $2$ przy $n=4$ podamy niżej.
\end{uwaga}

\begin{twierdzenie}[Odwrotne twierdzenie Kummera dla rozszerzeń cyklicznych]\label{aa:tw:kummer-converse}
Niech $F$ będzie ciałem, $n\geq2$ liczbą całkowitą
odwracalną w $F$, a $\zeta\in F$ pierwotnym
$n$-tym pierwiastkiem z jedynki. Jeśli $E/F$ jest
skończonym rozszerzeniem Galois, a $\operatorname{Gal}(E/F)$
jest grupą cykliczną rzędu $n$, to istnieją
$\alpha\in E^\times$ oraz generator $\sigma$ tej grupy takie, że
\[
E=F(\alpha),\qquad \alpha^n\in F,\qquad
\sigma(\alpha)=\zeta\alpha.
\]
\end{twierdzenie}
\begin{proof}
Wybierzmy dowolny generator $\sigma$ cyklicznej grupy
$G=\operatorname{Gal}(E/F)$. Rozszerzenie jest Galois,
więc $[E:F]=|G|=n$ i $E^G=F$.
Dla każdego $x\in E$ zdefiniujmy
\[
P(x)=\sum_{j=0}^{n-1}\zeta^{-j}\sigma^j(x).
\]
Jest to dobrze określone odwzorowanie $E\to E$, ponieważ
każdy składnik należy do $E$, a $\zeta$ jest niezerowe.
Automorfizmy $\sigma^0,\ldots,\sigma^{n-1}$ są parami różne.
Z lematu o niezależności homomorfizmów
\ref{aa:lem:dedekind-independent} oraz faktu, że pierwszy
współczynnik $\zeta^0=1$, wynika, że $P$ nie może być
identycznie zerowe. Wybierzmy $x\in E$ takie, że
$\alpha=P(x)\neq0$.

Aby sprawdzić, jak działa $\sigma$ na $\alpha$, przesuwamy
indeks w skończonej sumie. Ponieważ $\zeta\in F$ i $\sigma$
ustala $F$, mamy
\[
\begin{aligned}
\sigma(\alpha)
 &=\sum_{j=0}^{n-1}\zeta^{-j}\sigma^{j+1}(x)\\
 &=\zeta\sum_{k=1}^{n}\zeta^{-k}\sigma^k(x)
 =\zeta\left(\sum_{k=1}^{n-1}\zeta^{-k}\sigma^k(x)
       +\zeta^{-n}\sigma^n(x)\right)\\
 &=\zeta\left(\sum_{k=1}^{n-1}\zeta^{-k}\sigma^k(x)+x\right)
 =\zeta\alpha.
\end{aligned}
\]
W przedostatniej równości użyto $\zeta^n=1$
i $\sigma^n=\operatorname{id}_E$.
Z tej równości wynika
$\sigma(\alpha^n)=\sigma(\alpha)^n
=(\zeta\alpha)^n=\alpha^n$.
Ponieważ $\sigma$ generuje $G$, element $\alpha^n$
jest ustalony przez całe $G$, a więc $\alpha^n\in E^G=F$.

Pozostaje uzasadnić, że $\alpha$ generuje całe rozszerzenie.
Dla $0\leq j<n$ zachodzi
$\sigma^j(\alpha)=\zeta^j\alpha$.
Wartości te są parami różne, gdyż $\zeta$ ma rząd $n$
i $\alpha\neq0$.
Każdy automorfizm $F$-liniowy przenosi pierwiastek wielomianu
minimalnego $m_{\alpha,F}$ na inny jego pierwiastek.
Zatem $m_{\alpha,F}$ ma co najmniej $n$ różnych pierwiastków
w $E$, a więc $\deg m_{\alpha,F}\geq n$.
Z drugiej strony $F(\alpha)\subseteq E$, więc
$[F(\alpha):F]\leq[E:F]=n$.
Ponieważ $[F(\alpha):F]=\deg m_{\alpha,F}$,
obie nierówności dają $[F(\alpha):F]=n$ i $F(\alpha)=E$.
\end{proof}

\begin{przyklad}[Rozszerzenie kwadratowe]\label{aa:prz:kummer-square}
Dla $F=\mathbb Q$, $n=2$ oraz $a=2$ mamy
$\mu_2(F)=\{1,-1\}$, a liczba $2$ nie jest kwadratem
w $\mathbb Q$. Wielomian $X^2-2$ nie ma zatem pierwiastka
w $\mathbb Q$, jest nierozkładalny i
$E=\mathbb Q(\sqrt2)$ ma stopień $2$ nad $\mathbb Q$.
Jego oba pierwiastki $\sqrt2$ i $-\sqrt2$ należą do $E$;
automorfizmy Galois to identyczność oraz
$\sigma(u+v\sqrt2)=u-v\sqrt2$ ($u,v\in\mathbb Q$).
Grupa jest cykliczna rzędu $2$, a $\sigma(\sqrt2)=-\sqrt2$.
Jest to przypadek teorii Kummera z pierwotnym pierwiastkiem
$\zeta=-1$.
\end{przyklad}

\begin{przyklad}[Dlaczego potrzebujemy pierwiastków z jedynki?]\label{aa:prz:kummer-missing-roots}
Niech $\alpha=\sqrt[3]{2}$ będzie dodatnim rzeczywistym
pierwiastkiem wielomianu $X^3-2$, a
$\omega\in\mathbb C$ pierwotnym pierwiastkiem trzeciego stopnia
z jedynki, zatem $\omega^2+\omega+1=0$.
Wielomian $X^3-2$ jest nierozkładalny nad $\mathbb Q$
na mocy kryterium Eisensteina dla liczby pierwszej $2$,
więc $[\mathbb Q(\alpha):\mathbb Q]=3$.
Ciało $\mathbb Q(\alpha)$ zawiera się w $\mathbb R$,
a pozostałe pierwiastki $\omega\alpha,\omega^2\alpha$
są nierzeczywiste, zatem do niego nie należą.
Nie jest ono normalne nad $\mathbb Q$, a więc nie jest Galois.

Zastąpmy teraz ciało podstawowe przez $F'=\mathbb Q(\omega)$.
Mamy $[F':\mathbb Q]=2$ i $\mu_3(F')=\langle\omega\rangle$.
Gdyby $X^3-2$ był rozkładalny nad $F'$, to jako wielomian
stopnia $3$ miałby pierwiastek $\omega^j\alpha\in F'$
dla pewnego $j\in\{0,1,2\}$. Wówczas
$\alpha=\omega^{-j}(\omega^j\alpha)\in F'$, co przeczy
$[\mathbb Q(\alpha):\mathbb Q]=3$, gdyż stopień podciała
$\mathbb Q(\alpha)\subseteq F'$ musiałby dzielić
$[F':\mathbb Q]=2$.
Zatem $X^3-2$ jest nierozkładalny nad $F'$, a
$E=F'(\alpha)/F'$ jest cyklicznym rozszerzeniem Galois
stopnia $3$ z automorfizmem $\sigma(\alpha)=\omega\alpha$.
Jego działanie na pierwiastkach przedstawia
rysunek~\ref{aa:fig:kummer-cubic-cycle}.
\end{przyklad}
\begin{figure}[htbp]
\centering
\begin{tikzpicture}[>=Stealth, every node/.style={font=\small}]
  \node (a) at (0,1.6) {$\alpha$};
  \node (b) at (2.6,0) {$\omega\alpha$};
  \node (c) at (-2.6,0) {$\omega^2\alpha$};
  \draw[->,thick] (a) -- node[above right] {$\sigma$} (b);
  \draw[->,thick] (b) -- node[below] {$\sigma$} (c);
  \draw[->,thick] (c) -- node[above left] {$\sigma$} (a);
\end{tikzpicture}
\caption{Cykliczne działanie generatora grupy Galois na trzech pierwiastkach $X^3-2$ nad $\mathbb Q(\omega)$.}
\label{aa:fig:kummer-cubic-cycle}
\end{figure}

\begin{przyklad}[Dlaczego charakterystyka ma znaczenie?]\label{aa:prz:kummer-inseparable}
Niech $p$ będzie dowolną liczbą pierwszą, $F=\mathbb F_p(t)$
(ciało funkcji wymiernych w zmiennej $t$), a $\alpha$
niech spełnia $\alpha^p=t$. Wielomian $X^p-t$ ma
pochodną zerową. Wykażemy, że $\alpha\notin F$.
Gdyby $t=(f(t)/g(t))^p$ dla niezerowych wielomianów
$f,g\in\mathbb F_p[t]$, to porównanie wykładników czynnika
nierozkładalnego $t$ po obu stronach równości
$f(t)^p=tg(t)^p$ dałoby
$p\,v_t(f)=1+p\,v_t(g)$, co jest niemożliwe modulo $p$.
Tutaj $v_t(h)$ oznacza krotność czynnika $t$ w niezerowym
wielomianie $h$.

Ponieważ $t=\alpha^p$ jest przestępne nad $\mathbb F_p$,
także $\alpha$ jest przestępne nad $\mathbb F_p$. Mamy
$F(\alpha)=\mathbb F_p(\alpha)$ i
$F=\mathbb F_p(\alpha^p)$.
Elementy $1,\alpha,\ldots,\alpha^{p-1}$ są liniowo
niezależne nad $F$: po sprowadzeniu ewentualnej relacji
do wspólnego mianownika otrzymalibyśmy niezerowy
wielomian w $\alpha$ o wyrazach, których wykładniki
są różne modulo $p$, sprzecznie z przestępnością $\alpha$.
Z drugiej strony $\alpha^p=t\in F$, więc
wielomiany w $\alpha$ redukują się do kombinacji
$1,\ldots,\alpha^{p-1}$. Zbiór tych kombinacji jest
podpierścieniem $F[\alpha]$ o wymiarze $p$ nad $F$;
ponieważ jest dziedziną skończenie wymiarową nad ciałem,
każdy niezerowy element ma odwrotność: mnożenie przez
niego jest iniektywnym, a więc surjektywnym
$F$-liniowym endomorfizmem $F[\alpha]$.
Zatem $F[\alpha]=F(\alpha)$, a wskazane elementy
tworzą bazę tego ciała nad $F$. W konsekwencji
$[F(\alpha):F]=p$.

Jeśli $\tau$ jest dowolnym $F$-automorfizmem $F(\alpha)$,
to $\tau(\alpha)^p=t=\alpha^p$, zatem
$(\tau(\alpha)-\alpha)^p=0$. W ciele element
o zerowej $p$-tej potędze musi być zerowy,
a więc $\tau(\alpha)=\alpha$ i $\tau=\operatorname{id}$.
Rozszerzenie nie jest Galois mimo stopnia $p$.
Właśnie dlatego nie wolno usuwać założenia
$\operatorname{char}F\nmid n$.
\end{przyklad}

\begin{przyklad}[Złożony wykładnik — stopień rozszerzenia może być mniejszy]
Rozpatrzmy $F=\mathbb Q(i)$, $n=4$ oraz $a=4$.
Ciało $F$ zawiera wszystkie cztery pierwiastki z jedynki
$\{1,-1,i,-i\}$. Wybierając $\alpha=\sqrt2$, mamy
$\alpha^4=4$, ale $[F(\alpha):F]=2$, ponieważ
$\sqrt2\notin\mathbb Q(i)$, gdyż rzeczywiste elementy $\mathbb Q(i)$
są wymierne, a $\sqrt2\notin\mathbb Q$; ponadto $\sqrt2$
spełnia $X^2-2=0.

Co więcej, $4$ nie jest czwartą potęgą w $F$.
Jeśli bowiem $b=x+yi\in\mathbb Q(i)$ ($x,y\in\mathbb Q$)
i $b^4=4$, to $b^2=2$ lub $b^2=-2$.
Ponieważ $b^2=(x^2-y^2)+2xyi$ jest rzeczywiste,
$xy=0$. Jeśli $b^2=2$, to $x=0$ jest niemożliwe,
więc $y=0$ i $x^2=2$. Jeśli $b^2=-2$, to $y=0$
jest niemożliwe, więc $x=0$ i $y^2=2$;
żaden z tych warunków nie zachodzi dla liczb wymiernych.
Tak więc $a$ nie jest czwartą potęgą, a mimo to
$X^4-4=(X^2-2)(X^2+2)$ ma nad $F$ czynniki stopnia $2$.
Grupa $\operatorname{Gal}(F(\sqrt2)/F)$ jest cykliczna
rzędu $2$, a nie $4$. Wyjaśnia to potrzebę dodatkowych
założeń przy klasyfikacji rozszerzeń o zadanym stopniu.
\end{przyklad}

\section{Kiedy dwa generatory dają to samo rozszerzenie?}
Dla liczby pierwszej $p$ sytuacja jest szczególnie przejrzysta:
nie istnieją nietrywialne właściwe podgrupy grupy cyklicznej
rzędu $p$. Dlatego przy założeniach $\operatorname{char}F\neq p$
oraz $|\mu_p(F)|=p$ element niebędący $p$-tą potęgą
daje rozszerzenie dokładnie stopnia $p$.

\begin{twierdzenie}[Klasyfikacja cyklicznych rozszerzeń Kummera stopnia $p$]\label{aa:tw:kummer-prime-classification}
Niech $p$ będzie liczbą pierwszą różną od
$\operatorname{char}F$. Załóżmy, że $F$ zawiera pierwotny
$p$-ty pierwiastek z jedynki. Niech $a,b\in F^\times$
nie będą $p$-tymi potęgami w $F$, a
$\alpha,\beta\in\overline F$ spełniają
$\alpha^p=a$ oraz $\beta^p=b$.
Wówczas równoważne są warunki:
\begin{enumerate}[label=\textup{(\roman*)}]
\item $F(\alpha)=F(\beta)$ jako podciała $\overline F$;
\item istnieją $j\in\{1,\ldots,p-1\}$ oraz
$c\in F^\times$ takie, że $b=c^p a^j$.
\end{enumerate}
\end{twierdzenie}
\begin{proof}
\emph{Przygotowanie: oba rozszerzenia mają stopień $p$.}
Na mocy twierdzenia~\ref{aa:tw:kummer-radical-cyclic}
$F(\alpha)/F$ jest Galois o grupie, której rząd dzieli $p$.
Jeśli rząd byłby $1$, to $[F(\alpha):F]=1$,
a więc $\alpha\in F$ i $a=\alpha^p$ byłoby $p$-tą
potęgą w $F$, sprzecznie z założeniem. Ponieważ $p$ jest
pierwsze, musi zachodzić $[F(\alpha):F]=p$.
Ten sam argument daje $[F(\beta):F]=p$.

\emph{Kierunek (ii)$\Rightarrow$(i).}
Załóżmy $b=c^p a^j$ dla wskazanych $c,j$.
Element $d=\beta/(c\alpha^j)$ jest dobrze określony,
ponieważ $c\alpha^j\neq0$, a jego $p$-ta potęga wynosi
\[
d^p=\frac{\beta^p}{c^p(\alpha^p)^j}
 =\frac{b}{c^pa^j}=1.
\]
Stąd $d\in\mu_p(\overline F)=\mu_p(F)\subseteq F$.
Otrzymujemy $\beta=dc\alpha^j\in F(\alpha)$,
czyli $F(\beta)\subseteq F(\alpha)$.
Obie strony mają stopień $p$ nad $F$, dlatego z twierdzenia
o stopniach w wieży wynika równość tych ciał.

\emph{Kierunek (i)$\Rightarrow$(ii).}
Niech $E=F(\alpha)=F(\beta)$, a $\zeta\in F$ będzie
pierwotnym pierwiastkiem $p$-tego stopnia z jedynki.
Grupa $\operatorname{Gal}(E/F)$ jest cykliczna rzędu $p$.
Każdy jej automorfizm wysyła $\alpha$ na $u\alpha$
dla pewnego $u\in\mu_p(F)$. Ponieważ $\alpha\notin F$,
nie wszystkie automorfizmy ustalają $\alpha$;
możemy więc wybrać generator $\sigma$ taki, że
$\sigma(\alpha)=\zeta\alpha$:
wyjściowy iloraz $u=\zeta^k$ ma $k\not\equiv0\pmod p$,
a odpowiednia potęga automorfizmu, wyznaczona odwrotnością
$k$ modulo $p$, daje żądany generator.

Ponieważ $\beta^p=b\in F$, także
$\sigma(\beta)/\beta$ jest w $\mu_p(F)$.
Istnieje zatem $j\in\{0,\ldots,p-1\}$ spełniające
$\sigma(\beta)=\zeta^j\beta$.
Wartość $j=0$ jest niemożliwa: wtedy $\sigma(\beta)=\beta$,
a ponieważ $\sigma$ generuje grupę Galois,
$\beta\in E^{\langle\sigma\rangle}=F$,
co zaprzecza temu, że $b$ nie jest $p$-tą potęgą w $F$.
Zatem $1\leq j\leq p-1$.
Dla elementu $c=\beta/\alpha^j\in E^\times$ obliczamy
\[
\sigma(c)=\frac{\zeta^j\beta}{(\zeta\alpha)^j}
=c.
\]
Ponieważ ciało elementów stałych całej grupy Galois jest równe $F$,
mamy $c\in F^\times$. Podnosząc równość
$\beta=c\alpha^j$ do potęgi $p$, dostajemy
$b=c^pa^j$, jak należało wykazać.
\end{proof}

\begin{uwaga}[Iloczyn klas a iloczyn tensorowy]\label{aa:uw:kummer-products}
Niech $p$ będzie liczbą pierwszą, a $F$ dowolnym ciałem.
Grupa $F^\times/(F^\times)^p$ składa się z warstw
$[a]=a(F^\times)^p$ dla $a\in F^\times$;
mnożymy je według reguły $[a][b]=[ab]$.
Każda klasa spełnia $[a]^p=[1]$, więc ta grupa abelowa
jest w naturalny sposób przestrzenią liniową nad $\mathbb F_p$:
dodawanie wektorów jest mnożeniem klas, a mnożenie
przez skalar $j\in\mathbb F_p$ jest operacją
$[a]\mapsto[a]^j$. Działania są dobrze określone,
ponieważ zmiana przedstawiciela $a$ o $p$-tą potęgę
nie zmienia klasy $[a]^j$.
Nie jest to ani iloczyn prosty grup, ani iloczyn tensorowy modułów.

Przy założeniu $\operatorname{char}F\neq p$ i obecności
pierwotnego $p$-tego pierwiastka z jedynki
niezerowe jednowymiarowe podprzestrzenie
$F^\times/(F^\times)^p$ odpowiadają dokładnie cyklicznym
rozszerzeniom Galois stopnia $p$ wewnątrz $\overline F$.
Istotnie, każda taka prosta jest generowana przez klasę
$[a]\neq[1]$, a dołączając $\sqrt[p]{a}$ otrzymujemy
rozszerzenie stopnia $p$ na mocy poprzedniego twierdzenia.
Dwie klasy generują tę samą prostą dokładnie wtedy,
gdy $[b]=[a]^j$ dla pewnego $1\leq j\leq p-1$;
poprzednie twierdzenie mówi, że jest to równoważne
równości odpowiednich ciał. Każde cykliczne rozszerzenie
Galois stopnia $p$ pojawia się w ten sposób na mocy
twierdzenia~\ref{aa:tw:kummer-converse}.
\end{uwaga}

\begin{przyklad}[Trzy różne rozszerzenia kwadratowe]\label{aa:prz:kummer-three-quadratic}
Niech $F=\mathbb Q$ i $p=2$. Ponieważ $-1\in\mathbb Q$
jest pierwotnym pierwiastkiem drugiego stopnia z jedynki,
możemy stosować twierdzenie~\ref{aa:tw:kummer-prime-classification}.
Liczby $2$, $3$ i $6$ nie są kwadratami liczb wymiernych.
Istotnie, jeśli $r^2=m$ dla dodatniej liczby całkowitej
$m\in\{2,3,6\}$ i liczby wymiernej $r$, to po zapisaniu
$r=u/v$ w postaci nieskracalnej, gdzie $u\in\mathbb Z$,
$v\in\mathbb Z_{>0}$, z równości $u^2=mv^2$ wynika,
że wykładnik każdej liczby pierwszej w rozkładzie $m$
musiałby być parzysty. Dla każdej z trzech liczb jest
to fałszywe.

Rozszerzenia $\mathbb Q(\sqrt2)$, $\mathbb Q(\sqrt3)$
i $\mathbb Q(\sqrt6)$ są zatem kwadratowe i Galois.
Są parami różne: dla $p=2$ zbiór dopuszczalnych wykładników
$j\in\{1,\ldots,p-1\}$ składa się wyłącznie z $j=1$,
więc równość dwóch z tych ciał wymagałaby, aby iloraz
odpowiednich liczb $2,3,6$ był kwadratem w $\mathbb Q$.
Ilorazy te to $2/3$, $1/3$ i $1/2$ (lub ich odwrotności),
a każdy ma w rozkładzie na czynniki pierwsze wykładnik
nieparzysty. Warunek klasyfikacji nie jest spełniony.

Jednocześnie $\sqrt6=\sqrt2\sqrt3$ należy do
$\mathbb Q(\sqrt2,\sqrt3)$. Iloczyn klas $[2][3]=[6]$
w $\mathbb Q^\times/(\mathbb Q^\times)^2$ odpowiada
trzeciemu podciału kwadratowemu rozszerzenia
$\mathbb Q(\sqrt2,\sqrt3)/\mathbb Q$; nie jest to
iloczyn tensorowy tych ciał ani ich iloczyn prosty.
\end{przyklad}

\begin{przyklad}[Dwie różne klasy dające to samo rozszerzenie]\label{aa:prz:kummer-class-example}
Niech $F=\mathbb Q(\omega)$, gdzie $\omega$ jest pierwotnym
pierwiastkiem trzeciego stopnia z jedynki,
a $a=2$, $b=4=a^2$. W przykładzie
\ref{aa:prz:kummer-missing-roots} wykazaliśmy,
że $2$ nie jest sześcianem w $F$.
Gdyby $[2]=[4]$ w $F^\times/(F^\times)^3$,
to $4/2=2$ byłoby sześcianem w $F$,
co jest niemożliwe. Klasy są więc różne.

Niech $\alpha^3=2$, $\beta=\alpha^2$.
Wtedy $\beta^3=4$ i oczywiście $\beta\in F(\alpha)$.
Z drugiej strony $\beta^2=\alpha^4=2\alpha$,
więc $\alpha=\beta^2/2\in F(\beta)$.
Stąd $F(\alpha)=F(\beta)$, mimo że $[2]\neq[4]$.
Obie klasy należą do tej samej jednowymiarowej podprzestrzeni,
ponieważ $[4]=[2]^2$, a w $\mathbb F_3$ skalar $2$
jest niezerowy.
Jeżeli $\sigma(\alpha)=\omega\alpha$, to
$\sigma(\beta)=\omega^2\beta$:
różne generatory tego samego rozszerzenia mogą więc
transformować się przez różne potęgi pierwotnego
pierwiastka z jedynki.
\end{przyklad}
